% =============================================================================
% Nota sulla copertina – pagina a sinistra, di fronte alla pagina del copyright
% =============================================================================
\thispagestyle{empty}
\vspace*{\fill}

\begin{center}
	{\sffamily\large\bfseries\textls[120]{\MakeUppercase{Sulla copertina}}}\par
	\vspace{0.6em}
	{\color{modernred}\rule{2.5cm}{0.8pt}}
\end{center}

\vspace{1.2em}

\noindent
{\small\justifying
	La figura in copertina è una testa in silhouette, rivolta verso l'alto, attraversata da una colata di colori e di frammenti digitali che scende dall'alto e la riempie. Quella colata sono i nostri bias: le scorciatoie mentali ereditate in milioni di anni di evoluzione, oggi amplificate dal flusso incessante di informazioni in cui viviamo. I circuiti che si intravedono all'interno del cranio ricordano che quel flusso non è più soltanto biologico: attraversa anche le macchine che abbiamo costruito.

	\smallskip
	\noindent Siamo immersi in questo flusso, ed è proprio per questo che è così difficile vederlo. Chi è trascinato da una corrente ne conosce la direzione, non la forma. Un occhio non può guardare se stesso, e una mente che pensa dentro i propri bias non può osservarli dall'interno: li scambia per la realtà. Ci sembra di ragionare, mentre in gran parte stiamo scorrendo.

	\smallskip
	\noindent Per guardare noi stessi occorre allora uscire, anche solo per un attimo, dal flusso. Fermarsi sulla riva, alzare lo sguardo come fa la figura in copertina, e vedere da fuori ciò che da dentro non si vede. Non è facile e non riesce mai fino in fondo: è un esercizio, non una conquista. Ma è in quell'attimo di distanza che nasce la possibilità di scegliere, invece di essere scelti.

	\smallskip
	\noindent Questo libro vuole essere quell'attimo: una sosta sulla riva, per guardare con curiosità, e senza giudizio, la corrente che ci porta.\par}

\clearpage
